\section*{الفصل \ref{ch:b3:canonical-ensemble} --- الجمهور القانوني}

\begin{solution}{exo:b3:canonical-ensemble:1}
(أ) $\mathcal P_2/\mathcal P_1 = (g_2/g_1)\eu^{-\beta\epsilon}$.
(ب) $\beta\epsilon = 2.1/(8.62\times10^{-5} \times 2500) = 9.7$: أي
نسبة $3\eu^{-9.7} \approx \num{2e-4}$. (ج) يحمل اللهب
$\sim\num{e15}$ ذرة صوديوم في cm$^3$: فحتى $10^{-4}$ منها،
وهي تدور كل نانوثوانٍ قليلة، تصبّ $10^{19}$ فوتونًا أصفر
في الثانية --- فتعمي البصر. (د) من $3\eu^{-\beta\epsilon} = 0.1$:
$T = \epsilon/(k_{\text{B}}\ln30) \approx \qty{7200}{K}$ --- أي
غلاف ضوئي نجمي، لا لهبًا.
\end{solution}

\begin{solution}{exo:b3:canonical-ensemble:2}
(أ) $z = 1 + \eu^{-\beta\epsilon} + \eu^{-2\beta\epsilon}$. (ب)
$\langle E\rangle = \epsilon(\eu^{-\beta\epsilon} +
2\eu^{-2\beta\epsilon})/z$: أي $\to 0$ عند $T$ المنخفضة، و $\to \epsilon$
(المستوي المتوسط) عند المرتفعة. (ج) $\mathcal P_1 = 1/(\eu^{\beta
\epsilon} + 1 + \eu^{-\beta\epsilon})$ تتزايد رتيبًا مع
$T$، ولا تدنو من حدّها الأعلى $1/3$ إلا عند $T \to \infty$. (د)
أوزان بولتزمان لا تنقص إلا مع الطاقة عند $T > 0$؛ وقلب
الجمهرات يحتاج إلى درجات الحرارة السالبة في
\cref{ex:b3:microcanonical-ensemble:negative}، وهي بعيدة المنال على أيّ
خزان.
\end{solution}

\begin{solution}{exo:b3:canonical-ensemble:3}
(أ) عامل بولتزمان على $E_p = mgh$ عند $T$ منتظم. (ب) $h_0 =
k_{\text{B}}T/mg = \qty{8.4}{km}$. (ج) $\eu^{-8848/8440} \approx
0.35$: أي ثلث جوّ --- ولذلك يحمل قاصدو القمم الأكسجين.
(د) الجوّ الحقيقي ليس متساوي الحرارة: فدرجة الحرارة تهبط مع
الارتفاع، ومنه ينحرف المنحني عن أُسّي واحد.
\end{solution}

\begin{solution}{exo:b3:canonical-ensemble:4}
(أ) ثلاث حركات انسحابية: $\tfrac32 R = \qty{12.5}{J/(mol.K)}$،
موافقةً تمامًا. (ب) أضف دورانين: $\tfrac52 R =
\qty{20.8}{J/(mol.K)}$، كما قيس: فالاهتزاز مجمَّد. (ج)
المتنبَّأ به $\tfrac72 R = \qty{29.1}{J/(mol.K)}$؛ والقيمة المقيسة
$\approx26$ تبيّن أن الاهتزاز ذاب جزئيًّا لا غير ($\theta_{\text{اهتزاز}}
\approx \qty{3400}{K}$). (د) $3R = \qty{24.9}{J/(mol.K)}$:
أي ديولونغ--بتي، تطيعه أكثر الفلزات عند حرارة الغرفة و
يخفق فيه الألماس --- وهي حكاية التجمّد في
\cref{exo:b3:canonical-ensemble:6}.
\end{solution}

\begin{solution}{exo:b3:canonical-ensemble:5}
(أ) $z = \eu^{-\beta\hbar\omega/2}/(1 - \eu^{-\beta\hbar\omega})$.
(ب) $\langle E\rangle = -\partial_\beta\ln z$ يعطي الصورة
المذكورة مع $\langle n\rangle = 1/(\eu^{\beta\hbar\omega} - 1)$. (ج)
$C \to k_{\text{B}}$ عند $T$ المرتفعة؛ و $C \approx k_{\text{B}}(\beta
\hbar\omega)^2\eu^{-\beta\hbar\omega} \to 0$ عند المنخفضة. (د) ومع
$\hbar\omega$ طاقةً لكمّ ضوئي، يكون $\langle n\rangle$ هو
عدد الفوتونات الحرارية لكل نمط: فقانون بلانك على بعد فصل
واحد.
\end{solution}

\begin{solution}{exo:b3:canonical-ensemble:6}
(أ) عدديًّا، $C = \tfrac12 \times 3Nk_{\text{B}}$ قرب
$T/\theta_{\text{E}} \approx 0.34$. (ب) $\theta_{\text{E}}/T =
4.3$: ومنه $C/3Nk_{\text{B}} = 4.3^2\eu^{-4.3}/(1 - \eu^{-4.3})^2
\approx 0.26$ --- فالألماس عند حرارة الغرفة لا يكاد يبلغ ربع
السعة الحرارية الكلاسيكية: أي إن ``الشذوذ'' تجمّد كمومي
على مرأى من الجميع. (ج) $\theta_{\text{E}} = \qty{90}{K}$
تضع الرصاص في عمق النظام الكلاسيكي عند \qty{300}{K}. (د) الفرضية
المخفقة هي فرضية التواتر الواحد: فللأجسام الصلبة الحقيقية طيف من الأنماط
نزولًا إلى الصوت الطويل الموجة، وكمّاته الرخيصة تعطي قانون $T^3$
--- أي إصلاح ديباي، في \cref{ch:b3:photons-phonons}.
\end{solution}

\begin{solution}{exo:b3:canonical-ensemble:7}
(أ) $f(v) \propto v^2\eu^{-mv^2/2k_{\text{B}}T}$: ومنه $v_{\text{p}} =
\sqrt{2k_{\text{B}}T/m}$: أي \qty{413}{m/s} من أجل N$_2$، و
\qty{1540}{m/s} من أجل H$_2$. (ب) $mv_{\text{إفلات}}^2/2k_{\text{B}}T
\approx 730$ من أجل N$_2$، و $52$ من أجل H$_2$. (ج) $\eu^{-730}$ يعني
أبدًا؛ و $\eu^{-52} \sim \num{e-23}$ لكل زمن مكوث بطيء ---
لكن الجوّ العلوي الساخن ($\sim\qty{1000}{K}$) يليّن
أُسّ الهدروجين إلى $\sim15$: فينزف الهدروجين على المدى
الجيولوجي، ويبقى الآزوت. (د) سرعة الإفلات ودرجة حرارة
الغلاف الخارجي: فبئر المشتري البالغ \qty{60}{km/s} يجعل أُسّ الهدروجين
نفسه فلكيًّا --- فالعمالقة الغازية تحفظ ما تفقده العوالم الصغيرة الدافئة.
\end{solution}

\begin{solution}{exo:b3:canonical-ensemble:8}
(أ) $\partial_\beta^2\ln Z = \langle E^2\rangle - \langle
E\rangle^2$، و $C = \partial_T\langle E\rangle =
-k_{\text{B}}\beta^2\partial_\beta\langle E\rangle$. (ب) يعطي
الطرفان $Nk_{\text{B}}(\beta\epsilon)^2\eu^{\beta\epsilon}/(
\eu^{\beta\epsilon} + 1)^2$. (ج) $\langle E\rangle \propto N$، و
$\Delta E \propto \sqrt N$. (د) مقدار استجابة جملة
للتسخين يساوي مقدار ارتجاف طاقتها تلقائيًّا ---
فالاستجابة والتذبذب قراءتان للمشتق الثاني
نفسه، وهو نمط (التذبذب--التبديد) يتكرر
في الفيزياء كلها.
\end{solution}

\begin{solution}{exo:b3:canonical-ensemble:9}
(أ) $\ln2 = E_a\,\Delta T/k_{\text{B}}T^2$ مع $\Delta T =
\qty{10}{K}$ و $T = \qty{298}{K}$: ومنه $E_a =
0.693\,k_{\text{B}}T^2/10 \approx \qty{0.55}{eV}$. (ب) من
$298$ إلى \qty{278}{K}: أبطأ بعامل $\eu^{E_a(1/278 - 1/298)/k_{\text{B}}}
\approx 4.4$ --- ولذلك تحفظ الثلاجات الطعام. (ج)
من $373 \to \qty{366}{K}$: تهبط السرعة بمقدار $\approx1.4$: فتحتاج
بيضة الإحدى عشرة دقيقة إلى ربع ساعة. (د) التفاعلات
يكسبها الذيل الأُسّي للجزيئات فوق الحاجز: فأزح
درجة الحرارة قليلًا تنزح جمهرة الذيل
انزياحًا هائلًا --- ولا يكاد المتوسط يعني شيئًا.
\end{solution}

\begin{solution}{exo:b3:canonical-ensemble:10}
(أ) $f_{\text{منطوي}} = 1/(1 + \eu^{-\beta(\Delta E -
T\Delta S)})$ مع إدماج إنتروبيا الحالة المنشورة في
طاقتها الحرة. (ب) عند $T_{\text{m}} = \Delta E/\Delta S$ تتعادل
الطاقتان الحرتان: نصفًا بنصف. (ج) $T_{\text{m}} =
\num{4.8e-19}/\num{1.38e-21} = \qty{348}{K}$ (أي \qty{75}{\celsius})؛
والعرض $\delta T \sim k_{\text{B}}T_{\text{m}}^2/\Delta E \approx
\qty{3.5}{K}$. (د) التعاونية تضرب $\Delta E$ و
$\Delta S$ كليهما في عدد التماسات المنكسرة معًا، فتبقي
$T_{\text{m}}$ وتقلّص العرض $\propto 1/\Delta E$: فتنفصل خيوط
الحمض النوويّ في حدود كلفنين، وهو ما يتيح لآلة
PCR أن تدور نظيفًا بين ``المنصهر'' و``المُلدَّن''.
\end{solution}

\begin{solution}{exo:b3:canonical-ensemble:11}
(أ) $z = 2\cosh(\beta\mu B)$؛ و $M = N\mu\tanh(\beta\mu B) \approx
N\mu^2B/k_{\text{B}}T$ من أجل وسيط صغير: أي $1/T$ كوري. (ب)
$\mu_{\text{B}}B/k_{\text{B}}T$: أي $\num{2.2e-3}$ عند \qty{300}{K}؛ و
$0.67$ عند \qty{1}{K} --- أي من اللامبالاة إلى المحاذاة القوية. (ج)
تتوقف $S$ على $\mu B/k_{\text{B}}T$ وحده؛ وإنقاص $B$ عند
إنتروبيا ثابتة يجرّ $T$ نزولًا بالتناسب: $\qty{1}{K} \to
\qty{0.1}{K}$. (د) حين يبلغ $k_{\text{B}}T$ طاقة التفاعل بين
سبين وسبين لا تعود الإنتروبيا محكومة بالمجال:
$\mu_0\mu_{\text{B}}^2/4\pi a^3 \approx \qty{7e-26}{J} \sim
\qty{5}{mK}$ --- وهي الأرضية الكلاسيكية (وأما العزوم النووية، وهي أضعف بألف
مرة، فتدفعها إلى الميكروكلفنات).
\end{solution}

\begin{solution}{exo:b3:canonical-ensemble:12}
(أ) $z_{\text{دوران}} = \sum_J(2J+1)\eu^{-\beta BJ(J+1)}
\to \int(2J+1)\eu^{-\beta BJ(J+1)}\dd J = k_{\text{B}}T/B$: ومنه
$\langle E\rangle = k_{\text{B}}T$ و $C_{\text{دوران}} =
k_{\text{B}}$. (ب) من أجل H$_2$: $\theta_{\text{دوران}} = \qty{87}{K}$ ---
فتقع الدرجة في مجال المخبر؛ ومن أجل N$_2$: \qty{2.9}{K}،
ولا تتجمّد إلا قرب الهليوم السائل، فيظهر الآزوت دائمًا
$\tfrac52 R$. (ج) المصاطب: $\tfrac32 R$ (تحت $\sim\qty{50}{K}$)، و
$\tfrac52 R$ (من $\sim200$ إلى $\sim\qty{1000}{K}$)، صاعدةً
نحو $\tfrac72 R$ قرب $\theta_{\text{اهتزاز}} \approx
\qty{6000}{K}$. (د) قيم $J$ الفردية والزوجية تنتمي إلى نوعين نوويين
من السبين يتحول أحدهما إلى الآخر ببطء، فتتوقف $C_V$ للهدروجين البارد
على تاريخه الأورثو--بارا --- أي إحصاء نوويّ يدقّقه
مسعّر.
\end{solution}

\begin{solution}{pb:b3:canonical-ensemble:1}
\textbf{1.} $V = \tfrac43\pi r^3 = \qty{4.0e-20}{m^3}$: ومنه $m' = V
\Delta\rho = \qty{8.3e-18}{kg}$، والوزن $m'g = \qty{8.1e-17}{N}$.
\textbf{2.} $n(h) = n_0\,\eu^{-m'gh/k_{\text{B}}T}$: أي
القانون البارومتري، مصغَّرًا إلى شريحة مجهر.
\textbf{3.} $h_0 = k_{\text{B}}T/m'g = \num{4.04e-21}/\num{8.1e-17}
\approx \qty{50}{\micro m}$.
\textbf{4.} $h_0 \propto 1/r^3$: فتشتّت $10\%$ في نصف القطر يعني
تشتّت $30\%$ في ارتفاع المقياس --- والحبيبات متعددة الأحجام تمسح
السلّم فتجعله عصيدة. وقد جزّأ بيران شهورًا بالنبذ
المتكرر.
\textbf{5.} الصيغة نفسها، والكتلتان متباعدتان بعامل $10^{10}$: فجوّ
الحبيبات أضحل $10^{10}$ مرة --- فتنكمش الكيلومترات
إلى عشرات الميكرومترات، وهو بالضبط ما يجعله
قابلًا للمشاهدة كله تحت مجهر.
\textbf{6.} النسب المتتالية $0.55$ و $0.55$ و $0.57$: ثابتة
في حدود خطأ العدّ --- أي أُسّية. ومنه $h_0 =
\qty{30}{\micro m}/\ln(100/55) \approx \qty{50}{\micro m}$.
\textbf{7.} $k_{\text{B}} = m'g\,h_0/T = \num{8.1e-17} \times
\num{5.0e-5}/293 \approx \qty{1.4e-23}{J/K}$.
\textbf{8.} $N_{\text{A}} = R/k_{\text{B}} \approx
\num{6.0e23}\,\unit{mol^{-1}}$.
\textbf{9.} في حدود بضعة بالمئة هنا (بمعطيات مثالية)؛ أما
تجارب بيران الحقيقية فتبعثرت بمقدار $\pm15\%$ حول القيمة الحديثة
--- وهو أمر مذهل لحبيبات معدودة باليد، وحاسم تمامًا
في رتبة القدر.
\textbf{10.} $m_{\text{H}} = \qty{e-3}{kg/mol}/N_{\text{A}} =
\qty{1.7e-27}{kg}$.
\textbf{11.} الكيمياء العيانية تقدّم $R$؛ والمجهر الضوئي
والوزن يقدّمان $m'$؛ والعدّ يقدّم السلّم: فثلاثة
قياسات على طاولة تثلّث كتلة ذرة لا يستطيع أحد
رؤيتها.
\textbf{12.} لا شيء: فلم يستعمل الاستنتاج إلا ``جملة
تتبادل الطاقة مع خزان'' --- فعامل بولتزمان
أعمى عن الحجم، وهو بالضبط ما تحقق منه بيران.
\textbf{13.} $D = k_{\text{B}}T/6\pi\eta r =
\num{4.04e-21}/\num{4.0e-9} \approx \qty{e-12}{m^2/s}$.
\textbf{14.} $\sqrt{2Dt} = \sqrt{\num{1.2e-10}} \approx
\qty{11}{\micro m}$ في الدقيقة: قابل للقياس بيسر بعدسة
عينية مدرَّجة وصبر.
\textbf{15.} ظاهرتان مستقلتان، وصيغتان مستقلتان،
وعدد واحد: فتوافق السلّم الساكن والارتجاف المتحرك
لم يترك ثغرة للمصادفة --- إذ تنبّأت الفرضية الجزيئية
بكليهما.
\textbf{16.} $\sqrt{k_{\text{B}}T/m} = \sqrt{\num{4.04e-21}/
\num{4.8e-17}} \approx \qty{9}{mm/s}$ لكل مركبة.
\textbf{17.} تُقرع الحبيبة $10^{19}$ مرة في الثانية و
تنسى سرعتها في حدود نانومترات: فالطيران القذفي
غير قابل للمشاهدة، وما تراه العين هو تكامله ---
أي المسير العشوائي الانتشاري.
\textbf{18.} سحب ستوكس (أي لزوجة الموائع في كتاب السنة
الثانية) يقدّم $6\pi\eta r$؛ \omterm{thm:b3:canonical-ensemble:boltzmann}{والجمهور القانوني} يقدّم
$k_{\text{B}}T$: فالمقدار $D$ عند أينشتاين خارج قسمتهما، أي ميكانيك
وإحصاء في كسر واحد.
\textbf{19.} الخيال لا يُعدّ: فما إن صار $N_{\text{A}}$
نسبةَ عددين مقيسين، بعصيَّ خطأ، حتى صارت الذرات
موضوعات تجربة --- وقد أقرّ أوستفالد بذلك مطبوعًا سنة 1909.
\textbf{20.} انتثار زرقة السماء، والتحليل الكهربائي بشحنة
الإلكترون المقيسة، والهليوم المتراكم من الراديوم: أربع قنوات
فيزيائية لا صلة بينها تقاربت على $6\times10^{23}$ واحد، فجعلت
العدد خاصةً من خواصّ الطبيعة لا من خواصّ نظرية.
\textbf{21.} طاقة ثقالية تجذب الحبيبات نزولًا،
وإنتروبيا تشكيلية تنشرها صعودًا، متقايضتين بالسعر $T$:
فالسلّم \emph{هو} أصغرية $F = E - TS$.
\textbf{22.} $h_0 \propto 1/g$: فعند $10^5g$،
$h_0 = k_{\text{B}}T/(m'\times10^5g) = \num{4.04e-21}/
\num{9.8e-17} \approx \qty{40}{\micro m}$ من أجل البروتين ---
أي توازن الترسيب، وهو تجربة بيران تُجرى يوميًّا في
أقسام الكيمياء الحيوية.
\textbf{23.} من ``أطول من اللازم'' ($h_0 \gg$ عمق الحقل: فلا تدرّج
مرئي) إلى ``أرقّ من اللازم'' ($h_0 \lesssim r$): فأنصاف الأقطار الصالحة تمتد
تقريبًا من $0.1$ إلى $\qty{0.5}{\micro m}$ --- ولم يكن اختيار بيران
حظًّا بل تصميمًا.
\textbf{24.} وما دام $k_{\text{B}}$ مضبوطًا الآن بالتعريف، فإن
التجربة نفسها تعاير الحبيبات (حجمها أو كثافتها) أو
تدقّق الجهاز: فالفيزياء --- أي سلّم بولتزمان --- لم
تُمسّ؛ ولم يتحرك إلا أيّ مقدار يُعدّ مجهولًا.
\textbf{25.} جمهرة حبيبة قطرها \qty{0.2}{\micro m} تنتصف كل
$\sim\qty{35}{\micro m}$؛ ومقروءًا عبر $n_0\eu^{-m'gh/k_{\text{B}}
T}$، أعطى السلّم $N_{\text{A}} \approx \num{6e23}$:
أي عدد أفوغادرو معدودًا حبيبةً حبيبة --- وأُغلق الجدل
الذرّي على مسرح مجهر.
\end{solution}
